\chapter{Exercício 6}
\label{cha:exercise_6}

\section{Aplicação Web GitTasks}

O exercício 6 consiste na implementação de uma aplicação web denominada \textbf{GitTasks}, responsável por integrar dados provenientes de repositórios GitHub com a criação de tarefas no serviço Google Tasks. A aplicação combina autenticação OpenID Connect, controlo de acesso baseado em papéis e interação com APIs externas através de pedidos HTTP.

A solução foi implementada integralmente sem SDKs oficiais, recorrendo apenas a chamadas HTTP diretas (via \texttt{axios}), sessões armazenadas em memória e controlo de permissões com Casbin.

\section{Configuração e variáveis de ambiente}

Para executar a aplicação, é necessário criar um ficheiro \texttt{.env} contendo as variáveis utilizadas pelo servidor:

\begin{verbatim}
PORT=3001
SESSION_SECRET=chave_ultra_hiper_mega_secreta_123
HMAC_SECRET=outra_chave_ultra_hiper_mega_secreta_mas_com_hmac_456

GOOGLE_CLIENT_ID=
GOOGLE_CLIENT_SECRET=
OAUTH2_REDIRECT_URI=http://localhost:3001/auth/callback

GITHUB_CLIENT_ID=
GITHUB_CLIENT_SECRET=
GITHUB_REDIRECT_URI=http://localhost:3001/github/callback
\end{verbatim}

As variáveis \texttt{GOOGLE\_CLIENT\_ID} e \texttt{GOOGLE\_CLIENT\_SECRET} são obtidas ao criar um cliente OAuth 2.0 na Google Cloud Console.  
As variáveis \texttt{GITHUB\_CLIENT\_ID} e \texttt{GITHUB\_CLIENT\_SECRET} são geradas ao criar uma OAuth Application em \texttt{github.com/settings/developers}.  
As chaves de sessão (\texttt{SESSION\_SECRET} e \texttt{HMAC\_SECRET}) são utilizadas pelo servidor para assinar cookies e proteger conteúdos sensíveis armazenados na sessão.

\section{Autenticação de utilizadores}

A autenticação principal é realizada via Google, utilizando o protocolo \textit{OpenID Connect}. O fluxo seguido pela aplicação é o seguinte:

\begin{itemize}
    \item O utilizador inicia a autenticação selecionando o botão Google.
    \item É gerado um valor \textit{state} protegido com HMAC e enviado para o endpoint de autorização do Google.
    \item Após o utilizador autenticar com a Google, o servidor recebe um código de autorização.
    \item O código é trocado por \textit{access token}, \textit{refresh token} e \textit{id token}.
    \item A aplicação pede os dados do utilizador ao endpoint \texttt{userinfo}.
    \item A sessão do utilizador é criada e persistida em memória, incluindo email, nome, avatar, tokens Google e o respetivo papel.
\end{itemize}

A sessão é mantida com cookies HTTP-only, garantindo isolamento entre utilizadores.

\section{Autorização e papéis (RBAC)}

A aplicação utiliza controlo de acessos baseado em papéis através da biblioteca \textbf{Casbin}.  
Toda a informação de papéis e permissões reside no ficheiro de política (\texttt{policy.csv}) e não no código.

\subsection{Atribuição de papéis}

Após o login do utilizador, a aplicação executa:

\begin{verbatim}
const role = await enforcer.getRolesForUser(email)
\end{verbatim}

O papel é obtido diretamente da política Casbin e é armazenado na sessão.  
Caso o email não exista na política, a autenticação é recusada.

\subsection{Permissões}

As permissões são definidas exclusivamente na política Casbin. Em termos gerais:

\begin{itemize}
    \item \textbf{free}: pode consultar milestones no GitHub.
    \item \textbf{regular}: pode criar tasks numa lista predefinida no Google Tasks.
    \item \textbf{premium}: pode criar tasks em listas cujo nome é definido pelo utilizador.
\end{itemize}

O middleware \texttt{requirePermissions} consulta o enforcer para validar se o utilizador tem permissão para a operação pretendida.

\section{Integração com GitHub}

A aplicação comunica diretamente com a API do GitHub, utilizando pedidos HTTP autenticados quando necessário:

\begin{itemize}
    \item Listar os repositórios do utilizador.
    \item Consultar detalhes de um repositório.
    \item Obter milestones de repositórios públicos ou privados.
\end{itemize}

Não é utilizado qualquer SDK; todas as chamadas são construídas manualmente com \texttt{axios}.

\section{Integração com Google Tasks}

A criação de tarefas no Google Tasks segue as permissões definidas no RBAC:

\begin{itemize}
    \item O utilizador \textbf{regular} cria tasks numa lista pré-definida pela aplicação.
    \item O utilizador \textbf{premium} pode indicar o nome da lista onde pretende criar a task.
\end{itemize}

O servidor gere automaticamente a expiração dos tokens, recorrendo ao \textit{refresh token} para obter novos access tokens quando necessário.

\section{Sessões e armazenamento em memória}

Toda a informação relativa ao utilizador permanece em memória:

\begin{itemize}
    \item Dados Google (email, nome, avatar).
    \item Tokens Google e respetivas datas de expiração.
    \item Estado da autenticação GitHub.
    \item Papel do utilizador atribuído via Casbin.
\end{itemize}

As sessões são identificadas por cookies HTTP-only, impedindo acesso via JavaScript.

\section{Interface Web}

A interface inclui:

\begin{itemize}
    \item Botões para login Google e GitHub.
    \item Informação do utilizador autenticado.
    \item Indicação clara do papel ativo.
    \item Formulários para criação de tarefas e consulta de repositórios e milestones.
    \item Comportamento dinâmico controlado pelo papel do utilizador.
\end{itemize}

\section{Segurança e controlo de acesso}

\begin{itemize}
    \item Endpoints sensíveis usam o middleware \texttt{requireAuth}.
    \item As permissões são verificadas com \texttt{enforcer.enforce(user, object, action)}.
    \item Tentativas de acesso indevido devolvem códigos HTTP 401 ou 403.
\end{itemize}

\section{Resumo}

A aplicação GitTasks cumpre todos os requisitos do exercício:

\begin{itemize}
    \item Integra GitHub e Google Tasks de forma transparente.
    \item Implementa autenticação completa via OpenID Connect.
    \item Aplica controlo de acessos RBAC baseado em Casbin.
    \item Mantém sessões e estados em memória sem necessidade de base de dados.
    \item Usa apenas HTTP requests diretos, sem dependência de SDKs externos.
    \item Adapta automaticamente a interface às permissões do utilizador.
\end{itemize}

A solução final apresenta uma arquitetura robusta, segura, e alinhada com as práticas solicitadas no enunciado.
